% ==============================================================================
% comparative_analysis.tex
% ==============================================================================

\section{Comparative Analysis}
\label{sec:comparative_analysis}

This section presents a structured comparison of the surveyed methods across key dimensions: architectural design, input requirements, evaluation metrics, and reported performance.

\subsection{Evaluation Metrics}

The following metrics are commonly used to assess reconstruction quality:

\begin{itemize}
  \item \textbf{Peak Signal-to-Noise Ratio (PSNR)}: Measures pixel-level fidelity between the reconstructed and ground-truth CT volumes.
  \item \textbf{Structural Similarity Index (SSIM)}: Captures perceptual similarity by comparing luminance, contrast, and structural information.
  \item \textbf{Fréchet Inception Distance (FID)}: Evaluates the distributional quality of generated volumes (primarily used for GAN and diffusion methods).
  \item \textbf{Mean Absolute Error (MAE)}: Quantifies average voxel-wise absolute deviation.
\end{itemize}

\subsection{Cross-Method Comparison}

% TODO: Insert comprehensive comparison table.
% Table should include: Method, Year, Architecture Family, Input Views,
% Dataset, PSNR, SSIM, and any notable characteristics.

\begin{table}[h]
  \centering
  \caption{Comparison of representative few-view X-ray-to-CT reconstruction methods. (To be completed.)}
  \label{tab:comparison}
  \begin{tabular}{llcccc}
    \toprule
    \textbf{Method} & \textbf{Family} & \textbf{Views} & \textbf{Dataset} & \textbf{PSNR} & \textbf{SSIM} \\
    \midrule
    X2CT-GAN (2019)  & GAN        & 2 & LIDC-IDRI & --    & --   \\
    NAF (2022)       & Neural Field & 10 & Custom   & --    & --   \\
    DPS (2023)       & Diffusion  & varies & varies & --    & --   \\
    \bottomrule
  \end{tabular}
\end{table}

\subsection{Key Observations}

\begin{enumerate}
  \item GAN-based methods generally produce sharper images but may introduce anatomically implausible details.
  \item Transformer-based methods show improved global coherence at the cost of increased computational requirements.
  \item Neural field approaches excel in resolution flexibility but require per-scene optimisation.
  \item Diffusion-based methods achieve state-of-the-art perceptual quality but are the slowest at inference time.
\end{enumerate}

% TODO: Expand observations with quantitative evidence from the literature.
